\section{Remediación priorizada}
\label{sec:remediacion}

La tabla~\ref{tab:remediacion} ordena las acciones por impacto y esfuerzo. Las
cuatro primeras concentraban la mayor reducción de riesgo y son, en su mayoría,
cambios de configuración o de filtros compartidos. La mayoría de estas acciones se
aplicó después de la evaluación; el estado exacto y las pruebas de verificación se
presentan en la sección~\ref{sec:estado-remediacion}.

\begin{table}[htbp]
\centering
\caption{Plan de remediación priorizado.}
\label{tab:remediacion}
\small
\setlength{\tabcolsep}{3pt}
\begin{tabular}{@{}C{1.6cm}L{7.2cm}L{3.3cm}L{2.7cm}@{}}
\toprule
\textbf{Prioridad} & \textbf{Acción} & \textbf{Hallazgos} & \textbf{Esfuerzo}\\
\midrule
1 & Autenticar y autorizar pacientes y triage; validar JWT en el gateway. & H-02, H-03, H-09, H-23 & Medio\\
2 & Restringir \texttt{/admin/**} a administradores y separar la gestión de usuarios. & H-04 & Bajo\\
3 & Eliminar el secreto JWT por defecto y rotarlo. & H-01 & Bajo\\
4 & Mover \texttt{env} fuera del repositorio y rotar Gmail y Groq. & H-05, H-13 & Bajo\\
5 & Actualizar Tomcat, FreeMarker, Jackson y HttpComponents. & H-06, H-14 & Medio\\
6 & Rate limiting y bloqueo en login y 2FA. & H-12, H-15 & Medio\\
7 & Corregir XSS y mass assignment. & H-07, H-08 & Medio\\
8 & Autenticar y limitar el agente IA. & H-10, H-21 & Medio\\
9 & Restringir CORS y añadir cabeceras de seguridad. & H-11, H-18 & Bajo\\
10 & Endurecer infraestructura, logs y manejo de errores. & H-19, H-20, H-22 & Medio\\
\bottomrule
\end{tabular}
\end{table}

\subsection{Horizonte de ejecución}

\begin{description}
  \item[Corto plazo (0 a 2 semanas).] Aplicar las prioridades 1 a 4 y la 9.
  Son cambios acotados que eliminan la posibilidad de leer o escribir datos
  clínicos sin credenciales, forjar tokens administrativos y exponer secretos.
  \item[Mediano plazo (2 a 6 semanas).] Actualizar dependencias, implantar rate
  limiting, corregir XSS y mass assignment, autenticar el agente y endurecer la
  infraestructura. Añadir pruebas de regresión para cada hallazgo.
  \item[Largo plazo (6 semanas o más).] Incorporar threat modeling al ciclo de
  desarrollo, adoptar ASVS como lista de verificación, establecer una política de
  dependencias y ejecutar ZAP/SCA/secret scanning de forma continua.
\end{description}

\subsection{Verificación}

Cada corrección debe verificarse con la misma evidencia que permitió detectar el
fallo. La tabla~\ref{tab:verificacion} propone las pruebas de aceptación.

\begin{table}[htbp]
\centering
\caption{Pruebas de aceptación sugeridas para las correcciones principales.}
\label{tab:verificacion}
\small
\setlength{\tabcolsep}{3pt}
\begin{tabular}{@{}L{4.5cm}L{7.4cm}L{3.3cm}@{}}
\toprule
\textbf{Corrección} & \textbf{Prueba de aceptación} & \textbf{Resultado esperado}\\
\midrule
Autenticación de APIs & Repetir \texttt{smoke-sin-token} y los recorridos IDOR. & 401 en todos los endpoints clínicos.\\
Autorización administrativa & Token de \texttt{doctor\_triage} contra \texttt{/admin/**}. & 403; token admin: 200.\\
Secreto JWT & Intentar firmar con el valor por defecto y con firma alterada. & 401 en ambos casos.\\
Secretos & Ejecutar gitleaks en árbol e historial. & 0 hallazgos reales.\\
Dependencias & Ejecutar Trivy SCA y comparar contra el umbral. & 0 críticas y 0 altas sin excepción.\\
2FA & 20 intentos fallidos y luego el código correcto. & Bloqueo temporal; no aceptación.\\
XSS & Repetir el payload en pacientes y seguimiento con Playwright. & Sin diálogo; salida codificada.\\
Mass assignment & Enviar \texttt{idPersonal} ajeno e inexistente. & Campo ignorado; 400/404.\\
CORS & Origen \texttt{evil.local} y origen permitido. & Sin cabecera para el origen no permitido.\\
Agente IA & Llamar sin token y con token de rol limitado. & 401/403; herramientas acotadas.\\
\bottomrule
\end{tabular}
\end{table}

\subsection{Controles preventivos}

\begin{itemize}
  \item \textbf{CI/CD.} Ejecutar gitleaks, Semgrep, Trivy y ZAP baseline en cada
  \textit{merge request}; bloquear críticas y altas sin justificación.
  \item \textbf{Pruebas de seguridad.} Añadir pruebas MockMvc para 401/403,
  validación de JWT, mass assignment y codificación de salida.
  \item \textbf{Secretos.} Usar variables de entorno o un gestor de secretos;
  prohibir valores por defecto en \path{application.properties}.
  \item \textbf{Dependencias.} Fijar versiones, revisar CVEs y mantener un
  \textit{software bill of materials} (SBOM).
  \item \textbf{Observabilidad.} Auditar accesos a historias clínicas, cambios de
  rol y creación de usuarios administrativos; no registrar tokens ni códigos.
  \item \textbf{Diseño.} Aplicar denegación por defecto, mínimo privilegio y
  autorización a nivel de objeto en todos los servicios.
\end{itemize}

